\documentclass[12pt,a4paper]{article}
\usepackage[a4paper,top=15mm,bottom=15mm,left=16mm,right=16mm]{geometry}
\usepackage[T1]{fontenc}
\usepackage{amsmath}
\usepackage{amssymb}
\usepackage{array}

\pagestyle{empty}
\setlength{\parindent}{0pt}
\setlength{\parskip}{4pt}
\allowdisplaybreaks

% A topic from the revision list, with its "Remember" line from the
% revision guide.
\newcommand{\ptitle}[2]{%
  \par\addvspace{10pt}
  \noindent\rule{\linewidth}{0.8pt}\par\nobreak\vspace{1pt}
  \noindent{\large\bfseries #1}\par\nobreak\vspace{1pt}
  \noindent\textbf{Remember: #2}\par\nobreak\vspace{1pt}
  \noindent\rule{\linewidth}{0.4pt}\par\nobreak\vspace{4pt}}
\newcommand{\wq}[1]{\par\addvspace{6pt}\noindent\textbf{Q#1.}\ }

\begin{document}

\begin{center}
{\Large\bfseries Standard form}\\[3pt]
{\large Revision questions: worked solutions}
\end{center}

% ----------------------------------------------------------------
\ptitle{Writing large and small numbers in standard form}{One digit before the point, then count the jumps.}

\wq{1} One digit before the point: 7.25. The 7 is in the $10^4$ column, so it jumps 4 places to get back to the ones column.
\[ 72\,500=7.25\times10^4 \]

\wq{2} (a) $6000=6\times10^3$ (3 jumps)\qquad (b) $3\,080\,000=3.08\times10^6$ (6 jumps)

(c) $0.009=9\times10^{-3}$ (3 jumps, and a small number has a negative power)

(d) $0.000\,304=3.04\times10^{-4}$ (4 jumps)

\wq{3} (a) $3\times10^6=3\,000\,000$\qquad (b) $2.07\times10^4=20\,700$

(c) $5\times10^{-2}=0.05$\qquad (d) $6.1\times10^{-5}=0.000\,061$

A positive power gives a big number; a negative power gives a number less than 1.

\wq{4} The first number must be from 1 up to, but not including, 10.

$34\times10^5$: 34 is too big. $34=3.4\times10^1$, so $34\times10^5=3.4\times10^6$.

$0.7\times10^3$: 0.7 is too small. $0.7\times10^3=700=7\times10^2$.

\wq{5} $58\,000\,000=5.8\times10^7$\,km (7 jumps)

\wq{6} Compare the powers first: $10^{-6}$ is the smallest, then $10^{-5}$, $10^{-4}$, $10^{-3}$.

Smallest first: red blood cell ($8\times10^{-6}$), human hair ($7\times10^{-5}$), grain of sand ($5\times10^{-4}$), ant ($4\times10^{-3}$).

% ----------------------------------------------------------------
\ptitle{Multiplying and dividing in standard form}{Numbers with numbers, powers with powers.}

\wq{7} Multiplying: add the powers. Dividing: subtract the powers.

(a) $2\times4=8$,\ $10^3\times10^5=10^8$:\quad $8\times10^8$

(b) $5\times3=15$,\ $10^6\times10^{-2}=10^4$:\quad $15\times10^4=1.5\times10^5$

(c) $9\div3=3$,\ $10^8\div10^2=10^6$:\quad $3\times10^6$

(d) $2\div8=0.25$,\ $10^{-3}\div10^4=10^{-7}$:\quad $0.25\times10^{-7}=2.5\times10^{-8}$

\wq{8} $(1\times10^{-4})\times(5\times10^3)=5\times10^{-1}$\,m, which is 0.5\,m.

\wq{9} Distance $=$ speed $\times$ time:\quad $(3\times10^5)\times(2\times10^3)=6\times10^8$\,km

% ----------------------------------------------------------------
\ptitle{Adding and subtracting in standard form}{Can't add powers: change to ordinary numbers first.}

\wq{10} (a) $60\,000+3000=63\,000=6.3\times10^4$

(b) $8\,500\,000-400\,000=8\,100\,000=8.1\times10^6$

(c) $0.002+0.0006=0.0026=2.6\times10^{-3}$

(d) $0.07-0.005=0.065=6.5\times10^{-2}$

\wq{11} Ben added the numbers ($3+2=5$) and added the powers ($4+3=7$). Powers are only added when we multiply. Change to ordinary numbers first:
\[ 30\,000+2000=32\,000=3.2\times10^4 \]

\wq{12} $143\,000-12\,700=130\,300=1.303\times10^5$\,km

% ----------------------------------------------------------------
\ptitle{Standard form with a calculator}{Use the $\times10^x$ key; it already includes the `$\times\,10$'.}

\wq{13} (a) $302\,400=3.024\times10^5$

(b) $5\,625\,000\,000=5.625\times10^9$

(c) $800\,000\,000=8\times10^8$

(d) $\sqrt{49\,000\,000}=7000=7\times10^3$

Check (a) without the calculator: $6.3\times4.8\approx30$ and $10^7\times10^{-3}=10^4$, so the answer is about $30\times10^4=3\times10^5$.

\wq{14} The calculator works from left to right. With no brackets, Ava's keys work out $4\times10^3=4000$, then $4000\div2=2000$, then $2000\times10^2=200\,000$. It has multiplied by $10^2$ instead of dividing by it.

Ava should put brackets round each number, or use the $\times10^x$ key, which keeps each number together:
\[ (4\times10^3)\div(2\times10^2)=4000\div200=20 \]

% ----------------------------------------------------------------
\ptitle{Problems with standard form}{Right formula, numbers in, answer in standard form.}

\wq{15} Time $=\dfrac{\text{distance}}{\text{speed}}=\dfrac{3.84\times10^8}{3\times10^8}=1.28$ seconds

($3.84\div3=1.28$ and $10^8\div10^8=10^0=1$)

\wq{16} Distance $=$ speed $\times$ time $=(1.1\times10^4)\times(3.6\times10^3)=3.96\times10^7$\,m

\wq{17} $\dfrac{2\times10^{-2}}{8\times10^{-6}}=0.25\times10^4=2.5\times10^3=2500$ red blood cells

\wq{18} $\dfrac{6.8\times10^7}{2.4\times10^5}=2.833\ldots\times10^2=283.3\ldots$, which is 280 to 2 significant figures ($2.8\times10^2$).

\wq{19} $\dfrac{5.97\times10^{24}}{7.35\times10^{22}}=0.8122\ldots\times10^2=81.22\ldots$, so the Earth is about 81 times heavier.

% ----------------------------------------------------------------
\par\addvspace{10pt}
\noindent\rule{\linewidth}{0.8pt}\par\nobreak\vspace{1pt}
\noindent{\large\bfseries Problem solving}\par\nobreak\vspace{1pt}
\noindent\rule{\linewidth}{0.4pt}\par\nobreak\vspace{4pt}

\wq{20} Time in seconds: $\dfrac{1.2\times10^{14}}{3\times10^9}=0.4\times10^5=4\times10^4$ seconds.

There are $60\times60=3600$ seconds in an hour: $40\,000\div3600=11.1\ldots$, so 11.1 hours.

\wq{21} The numbers multiply to 6 and the powers add to 5. For example
\[ (2\times10^2)\times(3\times10^3)\qquad (1.5\times10^1)\times(4\times10^4)\qquad (6\times10^{-1})\times(1\times10^6) \]

\wq{22} $(2\times10^3)^3=2^3\times10^9=8\times10^9$, and $10^{10}=10\times10^9$. So $10^{10}$ is bigger.

\end{document}
